% ============================================
% 现代化思维导图模板
% 专业学术风格 - 适用于概念关系展示
% ============================================
% Note: Colors (mind-*) are defined in the main document preamble (survey.ini.tex)
%
% Key improvements:
% 1. Increased sibling angles to prevent overlap
% 2. Larger level distances for better spacing
% 3. Wider text widths for longer content
% 4. Better font sizing hierarchy

\begin{figure}[H]
    \centering
    \begin{adjustbox}{max width=\textwidth, max height=0.85\textheight, center}
 \tikzset{
     % 注释样式 - 现代化设计
     every annotation/.style = {
         draw=mind-annotation-border!60,
         line width=1.2pt,
         fill=mind-annotation,
         font=\small\sffamily,
         rounded corners=4pt,
     },
     % 节点通用样式
     every node/.style = {
         concept,
         circular drop shadow={shadow xshift=1pt, shadow yshift=-1pt, opacity=0.15},
     },
     % 根节点样式 - 深色突出
     root/.style = {
         concept color=mind-root,
         font=\normalsize\sffamily\bfseries,
         text width=8em,
         text=white,
         minimum size=3cm,
     },
     % 一级节点样式 - 增大间距
     level 1 concept/.append style = {
         font=\small\sffamily\bfseries,
         sibling angle=72,
         text width=6em,
         level distance=16em,
         inner sep=4pt,
         minimum size=2.2cm,
     },
     % 二级节点样式 - 增大间距
     level 2 concept/.append style = {
         font=\footnotesize\sffamily,
         sibling angle=45,
         level distance=12em,
         text width=5em,
         minimum size=1.8cm,
     },
     % 三级节点样式
     level 3 concept/.append style = {
         font=\scriptsize\sffamily,
         sibling angle=40,
         level distance=9em,
         text width=4.5em,
         minimum size=1.5cm,
     },
 }
 
 % 信息注释命令 - 优化样式
 \newcommand*{\info}[4][15]{
     \node[annotation, #3, scale=0.7, text width=#1 em, inner sep=3mm] at (#2){
         \list{$\bullet$}{
             \topsep=2pt\itemsep=1pt\parsep=0pt\parskip=0pt
             \labelwidth=8pt\leftmargin=10pt\itemindent=0pt
             \labelsep=3pt
         }
         #4
         \endlist
     };
 }
 
     \begin{tikzpicture}
         \path[mindmap, concept color=mind-root, text=white, grow cyclic]
 <tree_code>
     \end{tikzpicture}
 \end{adjustbox}
 
 \caption{<caption>}
 \label{fig:<label>}
 
 \end{figure}

% ============================================
% USAGE INSTRUCTIONS FOR CODE GENERATION:
% ============================================
%
% When generating mindmap diagrams, follow these guidelines:
%
% 1. ROOT NODE:
%    - Use short, concise text (max 3-4 words)
%    - Example: node[root] {Open Challenges\\in LLM\\Function Calling}
%
% 2. LEVEL 1 NODES (main branches):
%    - Maximum 5-6 branches to avoid crowding
%    - Use sibling angle=72 for 5 branches, 60 for 6 branches
%    - Keep text short (2-3 words max)
%    - Example: child[concept color=mind-l1-1] { node {Multi-turn\\Reliability} }
%
% 3. LEVEL 2 NODES (sub-branches):
%    - Maximum 3-4 per parent
%    - Use sibling angle=45 for good spacing
%    - Can use slightly longer text with line breaks
%    - Example: child { node {Performance\\degradation} }
%
% 4. TEXT FORMATTING:
%    - Use \\ for line breaks in long text
%    - Keep each line under 10 characters
%    - Avoid very long single words
%
% 5. COLOR SCHEME:
%    - mind-l1-1: Blue (for technical topics)
%    - mind-l1-2: Green (for methodology)
%    - mind-l1-3: Yellow (for challenges)
%    - mind-l1-4: Pink (for applications)
%    - mind-l1-5: Cyan (for future work)
%    - mind-l1-6: Orange (for evaluation)
%
% EXAMPLE - Challenge mindmap with 4 main branches:
%
% node[root] {Open Challenges\\in LLM\\Function Calling}
%     child[concept color=mind-l1-1, grow=150] { 
%         node {Multi-turn\\Reliability}
%         child { node {State\\maintenance} }
%         child { node {Error\\recovery} }
%     }
%     child[concept color=mind-l1-2, grow=60] { 
%         node {Dynamic\\Environment}
%         child { node {Tool\\discovery} }
%         child { node {Schema\\evolution} }
%     }
%     child[concept color=mind-l1-3, grow=-30] { 
%         node {Reasoning\\Quality}
%         child { node {Tool\\selection} }
%         child { node {Parameter\\inference} }
%     }
%     child[concept color=mind-l1-4, grow=-120] { 
%         node {Operational\\Stability}
%         child { node {Latency\\control} }
%         child { node {Cost\\efficiency} }
%     };
%
% IMPORTANT SPACING RULES:
% - For 4 branches: grow angles = 150, 60, -30, -120 (90° apart)
% - For 5 branches: grow angles = 162, 90, 18, -54, -126 (72° apart)
% - For 3 branches: grow angles = 150, 30, -90 (120° apart)